\section{Related Work}
This section organizes prior work on mobility-on-demand systems by MDP components: \textit{State}, \textit{Action}, \textit{Transition}, and \textit{Reward}. For each component we first describe the majority modeling practice found in the literature, then highlight notable or less common alternatives and specialized approaches.

\subsection{State}
Many studies represent the system state with aggregated, zone- or node-level summaries. Optimization and equilibrium analyses encode counts and flows (idle vehicles, waiting passengers, OD demand rates, travel times) to keep the state compact and amenable to mathematical programming or dynamic programming \citep{spatiotemporal_pricing,spatial_pricing,geometric_matching}. This macroscopic view enables tractable computation of platform pricing and fleet control and is common in prescriptive models that aim for analytical insight.

Agent-based and simulation-driven works adopt richer, mesoscopic or microscopic state descriptions that include vehicle-level status, per-request attributes, and network link states for more realistic evaluation \citep{hybrid_operations,modeling_the_competition,competition_and_cooperation}. Learning-based papers often use structured intermediate representations — for instance, graph-structured node features or zonal vectors — to balance expressiveness and scalability, supporting transfer and GNN-based policies \citep{graph_meta,learning_based,competitive_pricing,two_sided}.

\subsection{Action}
Optimization-focused studies commonly formulate centralized, continuous controls such as zone prices, commission rates, rebalancing flows, or fleet sizing; these formulations are chosen to fit convex or decomposable optimization methods \citep{spatiotemporal_pricing,geometric_matching,spatial_pricing}.

Operational and ABM literature model discrete, localized actions (vehicle repositioning, offer construction, accept/decline decisions) reflecting per-vehicle or per-operator autonomy \citep{competition_and_cooperation,modeling_the_competition,hybrid_operations}. MARL and decentralized RL studies frequently adopt hybrid formulations: high-level continuous controls (desired distribution shares, price coefficients) issued by learned policies that are post-processed by optimization/solvers into executable discrete decisions \citep{learning_based,graph_meta,competitive_pricing}. Purely decentralized per-vehicle action spaces (rebalance/pickup/do nothing) are used in tabular or localized RL approaches such as SAMoD \citep{samod}.

\subsection{Transition}
Analytic and equilibrium studies often employ stylized, tractable transition models (network-flow ODEs, matching-time approximations, or static equilibrium constraints) to derive provable properties and bounds \citep{spatiotemporal_pricing,geometric_matching,spatial_pricing}. These models sacrifice some operational realism for analytical tractability.

Empirical, experimental and RL-focused papers rely on simulator-driven transitions that integrate routing, matching (KM/Hungarian), and vehicle state machines; such mesoscopic simulators enable realistic training and evaluation but require calibration (OSRM/Dijkstra routing, insertion heuristics, ILP re-optimizations) \citep{competitive_pricing,hybrid_operations,modeling_the_competition,competition_and_cooperation,graph_meta}. Meta-RL and transfer works also assume simulated task distributions across cities and use simulators to expose distributional shifts during training \citep{graph_meta}.

\subsection{Reward}
Optimization and economic analyses typically specify explicit objective functions such as platform profit maximization or social welfare (consumer and producer surplus). These works directly optimize or analyze these objectives using calculus, duality, or numerical optimization \citep{spatiotemporal_pricing,geometric_matching,spatial_pricing}.

In simulation and RL contexts rewards vary. MARL papers commonly use operator-level cumulative profit (revenue minus costs) as the reward signal \citep{competitive_pricing,learning_based,graph_meta}, while decentralized vehicle-level RL sometimes uses sparse or surrogate rewards (e.g., reward for serving passengers) to promote service-focused behavior \citep{samod}. Two-sided frameworks adopt composite reward designs balancing platform profit, commission collection, and service-level metrics to capture multi-stakeholder objectives \citep{two_sided}. ABM/game studies typically evaluate KPIs (profit, pooling efficiency, waiting time) without formulating an RL reward \citep{competition_and_cooperation,modeling_the_competition}.

\subsection{Synthesis}
Most papers sit on a spectrum between tractability and realism: aggregated states and continuous platform actions enable optimization and theory, whereas fine-grained states, discrete vehicle actions, and simulator-based transitions support realistic policy evaluation and learning. Graph-based and hybrid abstractions bridge these ends by preserving spatial structure for learning while remaining scalable \citep{graph_meta,learning_based,competitive_pricing}. Choosing among these practices depends on the research goal: theoretical insight and guarantees versus operational fidelity and policy discovery.
